\section{The Chain of Whys and the Restless Heart}

A creature held in being from moment to moment does not merely exist; it reaches. Every deliberate human action carries a quiet \emph{because} beneath it, and every finite \emph{because} invites another \emph{why}. A physician stays awake through the night to save a patient; the patient endures surgery to recover health; health is cherished to return to work and family; work and family are prized for security, affection, usefulness, and joy. Press any ordinary sequence of human choices far enough and the answers grow thinner until the speaker reaches something desired for its own sake. What stands at the far end of every human \emph{why}? And once justice has established what a creature owes to the Creator in worship, truth, and obedience, does the moral life terminate in an endless ledger of obligations, or does the structure of human desire point toward a good that completes the creature without dissolving its creaturehood?

St.\ Thomas Aquinas opens his moral theology with this question of the last end. Before naming a single commandment or classifying a single sin, he examines why human willing cannot drift through an infinite chain of intermediate aims. In \emph{Summa Theologiae} I-II, q.\,1, a.\,4, he observes that ends stand in an essential order of dependence: remove what comes first in an essentially ordered series, and everything subordinated to it loses its power to act. Human action involves two such orders moving in opposite directions---the order of intention and the order of execution. In the order of execution, there must be a first concrete step, or deliberation would circle forever and nobody would lift a hand. In the order of intention, there must be a final end that first awakens desire, for secondary goods move the will only insofar as they participate in or lead toward what is desired simply:

\begin{quote}
Consequently, on neither side is it possible to go to infinity since if there were no last end, nothing would be desired, nor would any action have its term, nor would the intention of the agent be at rest; while if there is no first thing among those that are ordained to the end, none would begin to work at anything, and counsel would have no term, but would continue indefinitely.\endnote{St.\ Thomas Aquinas, \emph{Summa Theologiae} I-II, q.\,1, a.\,4, corpus (English Dominican translation, Benziger/Project Gutenberg eBook 17897). Aquinas explicitly notes that an \emph{accidental} infinity of ends remains possible (for instance, multiplying reflective acts of willing or adding indefinite incidental aims), whereas an \emph{essentially} ordered chain of ends must terminate in a first principle of intention.}
\end{quote}

The claim is easy to misunderstand if one imagines a person reciting the name of a supreme good before tying a shoe or laughing at a table. Aquinas raises that exact objection against himself in \emph{ST} I-II, q.\,1, a.\,6: surely jokes, diversions, and everyday gestures are not consciously aimed at an ultimate horizon. His reply is grounded in ordinary life. One need not hold the last end before one's conscious attention in every decision, any more than a traveler walking down a country road must picture the destination at every stride. What remains active across every step is the \emph{virtue of the first intention}---the underlying orientation toward complete human fulfillment.\endnote{Aquinas, \emph{ST} I-II, q.\,1, a.\,6, corpus and ad 1--3. Even an action done in jest is aimed at the jester's own relaxation and good, which is desired as a part or beginning of complete good.} Whatever a human being chooses is chosen under the aspect of good; and because every partial good is a beginning that points toward completion, the will cannot want a fragment without implicitly wanting the whole in which its appetite could finally rest.

Yet human lives obviously pull in incompatible directions. One person bends every waking hour toward fortune, another toward political influence, another toward quiet comfort, another toward scientific discovery or devotion. How can all human beings share a single last end when their actual pursuits collide so sharply? Aquinas resolves the apparent contradiction in \emph{ST} I-II, q.\,1, a.\,7 by separating two senses of the last end: its formal character (\emph{ratio ultimi finis}) and the concrete reality in which a person expects to find it. Considered formally, every human being desires the same thing---the completion of one's perfection, the fullness of good that leaves nothing further to crave, which tradition calls happiness or beatitude. Considered materially, human beings divide over where that fullness resides. Aquinas borrows an image from the palate: every healthy tongue finds sweetness pleasant, yet one person prefers the sweetness of wine and another the sweetness of honey, and the truest verdict belongs to the palate that is not diseased. Even those who turn away from God through sin do not abandon the desire for beatitude; they misplace it, demanding from a finite substitute what only an infinite good can yield.\endnote{Aquinas, \emph{ST} I-II, q.\,1, a.\,7, corpus and ad 1. St.\ Augustine had stated the same point in \emph{De Trinitate} XIII.3, cited by Aquinas in the \emph{sed contra} (the ``on the contrary'' opening of the article).}

In the closing article of Question 1 (\emph{ST} I-II, q.\,1, a.\,8), Aquinas sharpens what sets the rational creature apart within creation. Every creature---stone, cedar, hawk, and human being---has God as its ultimate end considered as the supreme reality toward which the universe tends. Yet irrational creatures attain that end only by reflecting a created likeness of God, simply by existing, living, or sensing. ``For man and other rational creatures attain to their last end by knowing and loving God: this is not possible to other creatures, which acquire their last end, in so far as they share in the Divine likeness, inasmuch as they are, or live, or even know.''\endnote{Aquinas, \emph{ST} I-II, q.\,1, a.\,8, corpus.} A tree fulfills its nature by growing and bearing seed; a human person can be satisfied only by attaining the supreme good directly through intellect and will.

\subsection*{Aquinas's Examination of Created Goods}

If the will seeks complete rest, what in the created order can supply it? In Question 2 of the Prima Secundae, Aquinas weighs every candidate human beings have elevated into a life's center: wealth, honor, fame, power, bodily health, pleasure, the soul's own virtues and powers, and finally the whole cosmos of created goods. Every item on the list is a real good in its proper rank; the failure arises when a subordinate good is asked to bear the weight of an ultimate end.

Aquinas begins with \emph{wealth} (\emph{ST} I-II, q.\,2, a.\,1). Following Aristotle, he divides riches into natural wealth (food, shelter, clothing, tools that remedy bodily need) and artificial wealth (money invented to facilitate exchange). Natural wealth cannot be the last end because it exists for the sake of human life, standing beneath the person it serves; artificial wealth stands one step further down, sought only to procure natural necessities. To the objection that the desire for money seems infinite and therefore resembles the desire for the supreme good, Aquinas replies that the two infinities run in opposite directions:

\begin{quote}
For the more perfectly the sovereign good is possessed, the more it is loved, and other things despised: because the more we possess it, the more we know it. Hence it is written (Ecclus.\ 24:29): ``They that eat me shall yet hunger.'' Whereas in the desire for wealth and for whatsoever temporal goods, the contrary is the case: for when we already possess them, we despise them, and seek others: which is the sense of Our Lord's words (John 4:13): ``Whosoever drinketh of this water,'' by which temporal goods are signified, ``shall thirst again.'' The reason of this is that we realize more their insufficiency when we possess them: and this very fact shows that they are imperfect, and the sovereign good does not consist therein.\endnote{Aquinas, \emph{ST} I-II, q.\,2, a.\,1, ad 3. On the contrast between natural and artificial wealth, see the corpus of the same article.}
\end{quote}

Next come \emph{honor}, \emph{fame}, and \emph{power} (\emph{ST} I-II, q.\,2, aa.\,2--4). Honor fails because it does not even reside in the person honored; it lives in the mind and gesture of the one who offers deference, serving as a sign attesting an excellence that must already exist before applause arrives. To labor for honor rather than for the good itself turns virtue into ambition.\endnote{Aquinas, \emph{ST} I-II, q.\,2, a.\,2, corpus and ad 1.} Fame or human glory fares worse still, for Aquinas contrasts human knowledge with divine knowledge: human minds do not create the goodness they praise; they react to it from outside and frequently misjudge it, bestowing renown on falsehood and revoking applause on a rumor. Only the glory a creature possesses with God rests on a knowledge that causes goodness and cannot be deceived.\endnote{Aquinas, \emph{ST} I-II, q.\,2, a.\,3, corpus and ad 2--3.} Power, in turn, has the character of a \emph{principle} of action rather than a final \emph{end}, and it remains morally ambivalent---open equally to noble rule and cruel abuse. Borrowing from Boethius, Aquinas notes that the ruler surrounded by armed guards may inspire fear while feeling greater fear himself; human power neither cures anxiety nor guarantees goodness, whereas in God alone power and goodness are identical.\endnote{Aquinas, \emph{ST} I-II, q.\,2, a.\,4, \emph{sed contra}, corpus, and ad 1. Aquinas adds four general reasons why none of these four external ``goods of fortune'' can constitute beatitude: each can coexist with moral evil, each leaves essential goods missing, each can bring harm to its possessor, and each depends largely on external chance rather than on the principles within human nature.}

Turning inward from fortune to the person, Aquinas examines \emph{bodily goods}---health, strength, beauty, longevity (\emph{ST} I-II, q.\,2, a.\,5). A ship's captain takes meticulous care of the hull and rigging, yet the captain's ultimate aim is not to keep the vessel floating in harbor forever; the ship exists for the voyage. Because a human being is not the supreme good, the preservation of bodily existence cannot be the last end of reason and will, and the body itself stands ordered to the life of the soul as matter to form and instrument to agent.\endnote{Aquinas, \emph{ST} I-II, q.\,2, a.\,5, corpus. He adds in the \emph{sed contra} that if beatitude consisted in bodily excellence, elephants would outstrip humanity in longevity, lions in strength, and stags in swiftness.}

\emph{Pleasure} (\emph{ST} I-II, q.\,2, a.\,6) presents a subtler temptation because it feels self-justifying; as Aristotle noted, nobody asks why someone wishes to be pleased. Aquinas distinguishes the essence of a thing from its proper accident, the concomitant property that flows from it. Delight is the appetite's rest in a fitting good possessed, hoped for, or remembered. Even the highest spiritual joy is a consequence of possessing the supreme good rather than the essence of beatitude itself. As for bodily pleasure, it arises from sensible goods bounded by matter and locked to the singular, whereas the rational soul---independent of any bodily organ in its intellectual operation---opens onto universal truth and therefore possesses a kind of openness to infinity. Measured against the range of mind and will, sensible pleasure is ``quite a trifle as compared with the good of the soul.''\endnote{Aquinas, \emph{ST} I-II, q.\,2, a.\,6, corpus and ad 1, 3.}

Even the noble refuge of Stoic self-sufficiency---placing beatitude in \emph{some good of the soul}, whether virtue, knowledge, or inner poise---cannot survive scrutiny (\emph{ST} I-II, q.\,2, a.\,7). Aquinas draws a distinction that governs the rest of the treatise: ``end'' means either the reality we seek to attain (as money is the miser's object) or our attainment and enjoyment of it (the possession of money). As an attainment, beatitude indeed exists in the soul. Yet what \emph{makes} the soul happy cannot be the soul itself or any habit within it. The soul is by nature in potency, moving from ignorance to understanding and from weakness to virtue; and every good inherent in a created soul is a participated, portioned good. The human will, however, is ordered to the universal good. Nothing confined inside a finite soul can fill a capacity designed for more than the self.

The eighth article brings this weighing of created goods to its cliff edge (\emph{ST} I-II, q.\,2, a.\,8). Could the entire created universe, or communion with the angels, satisfy the human heart? Aquinas answers no:

\begin{quote}
It is impossible for any created good to constitute man's happiness. For happiness is the perfect good, which lulls the appetite altogether; else it would not be the last end, if something yet remained to be desired. Now the object of the will, i.e.\ of man's appetite, is the universal good; just as the object of the intellect is the universal true. Hence it is evident that naught can lull man's will, save the universal good. This is to be found, not in any creature, but in God alone; because every creature has goodness by participation. Wherefore God alone can satisfy the will of man, according to the words of Ps.\ 102:5: ``Who satisfieth thy desire with good things.'' Therefore God alone constitutes man's happiness.\endnote{Aquinas, \emph{ST} I-II, q.\,2, a.\,8, corpus (Ps 102:5 DR; Ps 103:5 Heb.); see also ad 2--3, where he explains that the whole universe of creatures is itself ordained to God as its end, and that while created goods are not smaller than what a human being can possess \emph{inherently}, they remain finite and restricted beside the infinite good of which the rational creature is capable \emph{as object}.}
\end{quote}

\subsection*{The Restless Heart Before Its Maker}

What Aquinas establishes through analysis of appetite and causality, St.\ Augustine had confessed centuries earlier in the opening lines of the \emph{Confessions}. Writing as a mortal sinner who had tested pleasure, rhetorical fame, philosophical ambition, and human affection one by one, Augustine begins by praising the Creator who planted longing in clay:

\begin{quote}
Great art Thou, O Lord, and greatly to be praised; great is Thy power, and of Thy wisdom there is no end. And man, being a part of Thy creation, desires to praise Thee, man, who bears about with him his mortality, the witness of his sin, even the witness that Thou ``resistest the proud,''---yet man, this part of Thy creation, desires to praise Thee. Thou movest us to delight in praising Thee; for Thou hast formed us for Thyself, and our hearts are restless till they find rest in Thee.\endnote{St.\ Augustine, \emph{Confessions} I.1.1 (translated by J.\ G.\ Pilkington, in \emph{Nicene and Post-Nicene Fathers}, First Series, vol.\ 1, Buffalo, 1886).}
\end{quote}

The \emph{Catechism of the Catholic Church} gathers this Augustinian and Thomistic convergence succinctly (no.\ 1718): the desire for happiness is of divine origin, placed in the human heart so that every lesser beauty and every disappointed craving might draw the person back to the One who alone fulfills it.

Yet this very discovery opens a stark crisis. If no created good can still the human will, and if God alone is the universal good toward which every human \emph{why} obscurely reaches, how could a finite, embodied intellect ever lay hold of the infinite God? Knowing that finite things fall short does not by itself place the Creator within reach. At the summit of Aquinas's examination of created goods, natural reason stands before a threshold it can recognize from below but cannot cross by its own stride.
